% !PRINTED 589
\weberpage{589}
Es muss demnach, wenn $\alpha\gamma$ durch $\beta$ theilbar ist, auch $\gamma$
durch $\beta$ theilbar sein. Also haben wir folgenden Fundamentalsatz:

5. \gesperrt{Ein Product aus zwei oder mehr ganzen Zahlen
in $R(i)$ ist nur dann durch eine Primzahl in $R(i)$
theilbar, wenn wenigstens einer seiner Factoren
durch diese Primzahl theilbar ist.}

Macht man nun noch die Bemerkung, dass eine Primzahl nur
dann durch eine andere theilbar sein kann, wenn der Quotient
eine Einheit ist, wenn also beide associirt sind, und betrachtet
associirte Primzahlen als nicht wesentlich verschieden, so folgt:

6. \gesperrt{Eine ganze Zahl des Körpers $R(i)$ kann immer
und nur auf eine Art in ein Product von Primzahlen
 zerlegt werden.}

Denn sei $\alpha$ irgend eine ganze Zahl in $R(i)$. Ist $\alpha$ nicht
selbst eine Primzahl, so ist sie zerlegbar, etwa in $\gamma\beta$; die Normen
von $\gamma$ und von $\beta$ sind aber kleiner als die Norm von $\alpha$. Ist $\beta$
noch keine Primzahl, so ist es wieder zerlegbar, etwa in $\delta\varepsilon$; die
Normen dieser Factoren, die alle positive ganze Zahlen sind,
nehmen immer ab, und man muss also nothwendig einmal auf
einen Factor stossen, der eine Primzahl ist. Also ist jede Zahl $\alpha$
gewiss durch eine Primzahl $\pi$ theilbar. Ist demnach $\alpha = \pi\alpha'$,
so gilt von $\alpha'$ dasselbe und es ist etwa $\alpha' = \pi'\alpha'' \dots$ Da die
Normen von $\alpha, \alpha', \alpha'' \dots$ wieder alle abnehmen, so muss man
bei der Fortsetzung dieser Reihe schliesslich auf eine Einheit
stossen. Es ist daher $\alpha$ in eine endliche Anzahl von Primfactoren
zerlegbar, also
\[
\alpha = \pi\pi'\pi'' \dots,
\]
wobei ein zuletzt übrig bleibender Einheitsfactor mit einer der
Primzahlen $\pi$ vereinigt werden kann. Sind nun $\varkappa, \varkappa', \varkappa'' \dots$
gleichfalls Primzahlen und ist $\alpha = \varkappa\varkappa'\varkappa'' \dots$, so folgt aus
\[
\varkappa\varkappa'\varkappa'' \dots = \pi\pi'\pi'' \dots
\]
nach 5., dass eine der Primzahlen $\pi$, z. B. die erste durch $\varkappa$ theilbar,
also von $\varkappa$ nicht wesentlich verschieden ist. Hebt man
beiderseits mit $\varkappa = \pi$, so kann man denselben Schluss mit $\varkappa'$
wiederholen u. s. f. und findet also, dass jede der Primzahlen $\varkappa$
auch unter den $\pi$ vorkommen muss, und wenn unter den $\varkappa$ eine
Primzahl mehrmals vorkommt, so muss sie mindestens ebenso
oft unter den $\pi$ vorkommen. Derselbe Schluss lässt sich aber
auch umgekehrt machen, und daraus folgt, dass die Gesammtheit
%CONT>>
